\documentclass[12pt]{article}
\usepackage[utf8]{inputenc}
\usepackage[T1]{fontenc}
\usepackage[margin=1in]{geometry}
\usepackage{setspace}
\usepackage{xcolor}
\usepackage[colorlinks=true, urlcolor=blue!60!black, linkcolor=blue!60!black]{hyperref}
\setlength{\parindent}{2em}
\setlength{\parskip}{1.25em}
\renewcommand{\baselinestretch}{1.0}
\begin{document}
\begin{center}
{\LARGE \textit{Shaw v. Reno}, 509 U.S. 630 (1993)\par}
\vspace{0.8em}
{\large Case Brief --- Four-Page Edition\par}
\end{center}
\vspace{1em}

\section{Memory Jogger}

Voters challenged North Carolina's snake-like, majority-Black District 12, drawn after a Justice Department objection; a 5--4 Court held that a redistricting plan so bizarre that it is unexplainable on grounds other than race states an Equal Protection claim triggering strict scrutiny, even without any dilution of a group's voting strength (\href{https://supreme.justia.com/cases/federal/us/509/630/}{509 U.S.\ 630}).

\section{Detailed Case Facts}

After the 1990 census, North Carolina, whose voting-age population was about 78\% white and 20\% Black, gained a twelfth House seat. Because 40 of its counties were covered by \S\ 5 of the Voting Rights Act (VRA), it submitted a plan with one majority-Black district for preclearance. The Attorney General objected that a second majority-minority district could have been drawn in the south-central to southeastern region.

Rather than seek a declaratory judgment, the General Assembly created two majority-Black districts: a hook-shaped District 1, likened to a ``Rorschach ink-blot test,'' and District 12, which snaked about 160 miles along Interstate 85, often no wider than the highway, splitting 10 counties to link dispersed Black neighborhoods. One legislator quipped that driving down I-85 with both car doors open would hit most of its residents. The Attorney General precleared the plan.

\section{Procedural History}

Five Durham County residents led by Ruth Shaw sued state and federal officials, alleging not vote dilution but that racially segregating voters into bizarre districts violated their right to a ``color-blind'' electoral process. A three-judge court dismissed the federal defendants and, 2--1, dismissed the claim against the state defendants, reading \textit{UJO v.\ Carey} (1977) to require discriminatory purpose and proportional underrepresentation. The Court noted probable jurisdiction, heard argument on April 20, 1993, and reversed on June 28, 1993.

\section{Judicial Votes}

Reversed and remanded, \textbf{5--4}.
\begin{itemize}
\item \textbf{Majority:} Justice O'Connor, joined by Chief Justice Rehnquist and Justices Scalia, Kennedy, and Thomas.
\item \textbf{Dissenting:} Justice White, joined by Justices Blackmun and Stevens; Justices Blackmun, Stevens, and Souter separately.
\end{itemize}

\section{Holding}

Appellants stated an analytically distinct Equal Protection claim by alleging that the plan was so irrational on its face that it could be understood only as an effort to segregate voters by race without sufficient justification. If that allegation stood uncontradicted on remand, the district court had to apply strict scrutiny, requiring the State to show narrow tailoring to a compelling interest.

\section{Analysis of Opinions}

\textbf{Majority (O'Connor).} Racial classifications are ``by their very nature odious to a free people'' and receive strict scrutiny whether called benign or remedial (\textit{Croson}). The rule reaches facially neutral laws ``unexplainable on grounds other than race'' (\textit{Arlington Heights}), including tortured boundaries (\textit{Gomillion}). Legislatures always know race, but disregard of compactness, contiguity, and political subdivisions is objective evidence of racial sorting. Racial sorting itself causes harm: it reinforces stereotypes that members of a race think alike, exacerbates bloc voting, and sends a ``pernicious'' message to representatives. \textit{UJO} involved vote dilution in a plan that followed traditional principles, and the relaxed political-gerrymandering standard of \textit{Davis v.\ Bandemer} does not apply to race. Section 5 is not ``carte blanche'' for racial gerrymandering beyond what is reasonably necessary; whether \S\ 2 or remedying past discrimination could justify District 12 was left for remand.

\textbf{White (dissenting).} Whites remained a majority in 10 of 12 districts (83\% of seats versus 79\% of voting-age population), so there was no cognizable injury. \textit{UJO}, \textit{Gomillion}, and \textit{Wright} require purpose and dilutive effect; focusing on ``superficialities of shape and size'' invents an unworkable claim.

\textbf{Blackmun, Stevens, and Souter (dissenting).} Blackmun noted the irony of striking a plan that sent North Carolina's first Black representatives since Reconstruction. Stevens argued that shape only evidences purpose, and that facilitating representation for a minority does not violate impartial governance. Souter stressed that race consciousness is unavoidable under the VRA and that district assignment denies no one the vote.

\textbf{Later developments.} \textit{Miller v.\ Johnson} (1995) made race's ``predominant'' role, not bizarreness, the trigger, and \textit{Shaw v.\ Hunt} (1996) struck District 12. \textit{Cooper v.\ Harris} (2017) required plaintiffs to ``disentangle race from politics''; \textit{Alexander v.\ South Carolina NAACP} (2024) presumed legislative good faith; and \textit{Louisiana v.\ Callais} (2026) held that \S\ 2 compliance is compelling only when \S\ 2 is read to enforce the Fifteenth Amendment's ban on intentional discrimination.

\section{Examples: Future Applications}

\begin{enumerate}
\item \textbf{Same-Side (Invalid): Barbell-Shaped Hispanic District.} A 210-mile district built to a ``50\%+ HVAP floor'' without a strong inference of intentional discrimination fails strict scrutiny under \textit{Callais}.
\item \textbf{Same-Side (Strict Scrutiny): Regular-Looking District With a Racial Target.} Moving tens of thousands of voters to hit a 50\% target triggers strict scrutiny even without a strange outline, as in \textit{Cooper v.\ Harris}.
\item \textbf{Opposite-Side (Claim Fails): Partisan Line-Drawing.} Moving Democratic Hispanic voters to protect an incumbent, with no alternative map, fails under \textit{Alexander} and \textit{Rucho}.
\item \textbf{Opposite-Side (Deferential Review): Race-Aware but Compact.} A compact, 55\% Black district explained by traditional criteria stands despite awareness of race.
\item \textbf{Fence-Sitter: Court-Ordered Remedial District.} A district drawn after a \S\ 2 finding faces strict scrutiny; whether the finding meets \textit{Callais}'s intent standard and the district is narrowly tailored is open.
\end{enumerate}

\section{Critique}

\textbf{Standing and expressive harm.} \textit{Shaw} plaintiffs suffer no vote dilution. Critics argue that the ``expressive harm'' theory turns aesthetic distaste into constitutional injury; defenders respond that state-sponsored racial sorting itself violates individual dignity.

\textbf{The race--partisanship paradox.} Together, \textit{Shaw}, \textit{Rucho}, \textit{Alexander}, and \textit{Callais} create an asymmetry: sorting Black voters because they are Black triggers near-fatal strict scrutiny, while sorting the same voters because they vote Democratic is permissible and nonjusticiable in federal court.

\section{Key Quotations}

\begin{quote}
``Racial gerrymandering, even for remedial purposes, may balkanize us into competing racial factions; it threatens to carry us further from the goal of a political system in which race no longer matters---a goal that the Fourteenth and Fifteenth Amendments embody, and to which the Nation continues to aspire.''
\end{quote}
\noindent --- Justice O'Connor, 509 U.S.\ at 657.

\begin{quote}
``When a district obviously is created solely to effectuate the perceived common interests of one racial group, elected officials are more likely to believe that their primary obligation is to represent only the members of that group, rather than their constituency as a whole. This is altogether antithetical to our system of representative democracy.''
\end{quote}
\noindent --- Justice O'Connor, 509 U.S.\ at 648.

\begin{quote}
``Put differently, we believe that reapportionment is one area in which appearances do matter. A reapportionment plan that includes in one district individuals who belong to the same race, but who are otherwise widely separated by geographical and political boundaries, and who may have little in common with one another but the color of their skin, bears an uncomfortable resemblance to political apartheid.''
\end{quote}
\noindent --- Justice O'Connor, 509 U.S.\ at 647.

\end{document}
